\documentclass[14pt]{extarticle}
\usepackage{fontspec}
\usepackage{polyglossia}
\usepackage{amsmath}
\usepackage{amssymb}
\usepackage{xcolor}
\usepackage{fancyhdr}
\usepackage{graphicx}
\usepackage{listings}
\usepackage[most]{tcolorbox}
% \usepackage{setspace}
\usepackage{pifont}
\usepackage{enumitem}
\usepackage{titlesec}
\usepackage[bottom]{footmisc}
\usepackage{titling}
\usepackage{minted}
\usepackage{etoolbox}
\usepackage{array}
\usepackage{longtable} % page-breakable tables (must load before polyglossia/bidi RTL setup)
\usepackage{extsizes}
\usepackage{float}
\usepackage{booktabs}
\usepackage{hyperref}
\usepackage[onehalfspacing]{setspace}
\newif\ifwhiteblack
\ifdefined\setwhiteblack
  \setwhiteblack
\fi

\ifwhiteblack
  \usepackage[a4paper, landscape, top=1.25cm, bottom=1.25cm, left=1.25cm, right=1.25cm, headheight=1.5cm, headsep=0.5cm]{geometry}
\else
  \usepackage[a4paper, top=2.7cm, bottom=2.7cm, left=2.5cm, right=2.5cm, headheight=1.5cm, headsep=0.5cm]{geometry}
\fi
\usepackage{tikz}
\usepackage{datetime2}

\usetikzlibrary{automata, positioning, arrows.meta, arrows}
\providecommand{\LTcaptype}{table}

\newfontfamily\emoji{Segoe UI Emoji}

\pagestyle{fancy}

\setmainlanguage[numerals=western]{arabic}
\setotherlanguage{english}
\newfontfamily\arabicfont[Script=Arabic]{Calibri}
\newfontfamily\arabicfonttt[Script=Arabic]{Courier New}



\ifwhiteblack
\lstset{
  language=[Sharp]C,
  numbers=left,
  stepnumber=1,
  numbersep=8pt,
  frame=single,
  basicstyle=\ttfamily\small,
  keywordstyle=\bfseries,
  stringstyle=\ttfamily,
  commentstyle=\itshape
}
\else
\lstset{
  language=[Sharp]C,
  numbers=left,
  stepnumber=1,
  numbersep=8pt,
  frame=single,
  basicstyle=\ttfamily\small,
  keywordstyle=\color{blue},
  stringstyle=\color{red},
  commentstyle=\color{green!50!black}
}
\fi

\newif\ifdetailed
\ifdefined\setdetailed
  \setdetailed
\fi

\newif\ifwithsols
\ifdefined\setwithsols
  \setwithsols
\fi

\ifwhiteblack
  \AtBeginDocument{
    \LARGE
    \usemintedstyle{bw}
  }
\fi

% Default minted look: compact code that wraps long lines.
% Not applied in black-and-white landscape mode, which already uses \LARGE.
\ifwhiteblack\else
  \setminted{fontsize=\small, breaklines=true}
\fi

% unified tcolorboxes for articles
\ifwhiteblack
\tcbset{colback=white, colframe=black, fonttitle=\bfseries, boxrule=1pt}
\newtcolorbox{boxDef}[1][]{title={{\emoji📘} تعريف\ifx\\#1\\\else ~#1\fi :}}
\newtcolorbox{boxExercise}[1][]{title={{\emoji🧩} تمرين\ifx\\#1\\\else ~#1\fi :}}
\newtcolorbox{boxExample}[1][]{title={{\emoji📝} مثال\ifx\\#1\\\else ~#1\fi :}}
\newtcolorbox{boxNote}[1][]{title={{\emoji✨} ملاحظة\ifx\\#1\\\else ~#1\fi :}}
\newtcolorbox{boxAttention}[1][]{title={{\emoji🔔} تنبيه\ifx\\#1\\\else ~#1\fi :}}
\newtcolorbox{boxWarning}[1][]{title={{\emoji⚡} ملاحظة هامة\ifx\\#1\\\else ~#1\fi :}}
\newtcolorbox{boxSolution}[1][]{title={{\emoji✅} حل\ifx\\#1\\\else ~#1\fi :}}
\newtcolorbox{boxSymbol}[1][]{title={{\emoji🔣} رمز\ifx\\#1\\\else ~#1\fi :}}
\newtcolorbox{boxHint}[1][]{title={{\emoji💡} تلميح\ifx\\#1\\\else ~#1\fi :}}
\else
\tcbset{colback=white, colframe=black, fonttitle=\bfseries, boxrule=0.8pt}
\newtcolorbox{boxDef}[1][]{colback=blue!5!white,colframe=blue!75!black,
  title={{\emoji📘} تعريف\ifx\\#1\\\else ~#1\fi :}}
\newtcolorbox{boxExercise}[1][]{colback=cyan!5!white,colframe=cyan!70!black,
  title={{\emoji🧩} تمرين\ifx\\#1\\\else ~#1\fi :}}
\newtcolorbox{boxExample}[1][]{colback=yellow!5!white,colframe=orange!90!black,
  title={{\emoji📝} مثال\ifx\\#1\\\else ~#1\fi :}}
\newtcolorbox{boxNote}[1][]{colback=gray!10!white,colframe=black,
  title={{\emoji✨} ملاحظة\ifx\\#1\\\else ~#1\fi :}}
\newtcolorbox{boxAttention}[1][]{colback=magenta!10!white,colframe=magenta!80!black,
  title={{\emoji🔔} تنبيه\ifx\\#1\\\else ~#1\fi :}}
\newtcolorbox{boxWarning}[1][]{colback=red!5!white,colframe=red!75!black,
  title={{\emoji⚡} ملاحظة هامة\ifx\\#1\\\else ~#1\fi :}}
\newtcolorbox{boxSolution}[1][]{colback=green!5!white,colframe=green!60!black,
  title={{\emoji✅} حل\ifx\\#1\\\else ~#1\fi :}}
\newtcolorbox{boxSymbol}[1][]{colback=purple!5!white,colframe=purple!70!black,
  title={{\emoji🔣} رمز\ifx\\#1\\\else ~#1\fi :}}
\newtcolorbox{boxHint}[1][]{colback=teal!5!white,colframe=teal!60!black,
  title={{\emoji💡} تلميح\ifx\\#1\\\else ~#1\fi :}}
\fi


\ifwhiteblack
\tcbset{simplecode/.style={ colback=white, colframe=black, boxrule=1pt, arc=2pt, left=4pt,right=4pt,top=4pt,bottom=4pt}}
\else
\tcbset{simplecode/.style={ colback=gray!5, colframe=black!50, boxrule=0.4pt, arc=2pt, left=4pt,right=4pt,top=4pt,bottom=4pt}}
\fi
% Compact code boxes: tighter line spacing (the preamble uses onehalfspacing) and small vertical gaps.
% Must come AFTER the simplecode definition above.
\tcbset{simplecode/.append style={fontupper=\setstretch{0.95}, top=4pt, bottom=4pt, before skip=6pt, after skip=6pt}}
\newenvironment{boxCode}{\begin{tcolorbox}[simplecode]}{\end{tcolorbox}}

% Code cell used inside tables (keeps code left-to-right inside RTL tables).
% Use inside a fixed-width p{...} column: \begin{pycell}\begin{minted}...\end{minted}\end{pycell}
\newtcolorbox{pycell}{frame empty,colback=white,boxsep=0pt,left=2pt,right=2pt,top=2pt,bottom=2pt,width=\linewidth}

\newcolumntype{C}[1]{>{\centering\arraybackslash}p{#1}}

\makeatletter
\DTMnewdatestyle{dotted}{%
  \renewcommand{\DTMdisplaydate}[4]{\two@digits{##3}.\two@digits{##2}.##1}%
  \renewcommand{\DTMDisplaydate}{\DTMdisplaydate}%
}
\makeatother
\DTMsetdatestyle{dotted}

% redefine spaces after titles
\makeatletter
\renewcommand{\@maketitle}{%
  \begin{center}
    {\huge \bfseries \@title \par}%
    \vskip 0.2em % space between title and author
    {\large \@author \par}%
    % \vskip 0.2em % space between author and date
    % {\normalsize \@date \par}%
    {\small \hfill \begin{english}\DTMtoday\end{english} \par}%
  \end{center}
}
\makeatother

\fancyhf{} % clear default
\fancypagestyle{plain}{
  \fancyhf{}
  \fancyhead[L]{مدرسة التسامح الشاملة}
  % \fancyhead[L]{\includegraphics[height=1cm]{../../../images/logoTasamoh.png}}
  \fancyhead[R]{الأستاذ محمود اغبارية}
  \fancyfoot[C]{\thepage}
}

\fancyhead[L]{مدرسة التسامح الشاملة}
\fancyhead[R]{الأستاذ محمود اغبارية}
\fancyfoot[C]{\thepage}

\setcounter{tocdepth}{3} % only section subsection and subsubsection in TOC

\newcommand{\importpdfpage}[2]{%
    \thispagestyle{empty}
    \newgeometry{left=0cm,right=0cm,top=0cm,bottom=0cm}
    \includegraphics[page=#2,width=0.95\paperwidth,height=0.95\paperheight,keepaspectratio]{#1}
    \restoregeometry
}

\newcommand{\insertFullImg}[1]{%
    \noindent
    \makebox[\textwidth][c]{
        \includegraphics[width=0.9\paperwidth,keepaspectratio]{#1}%
    }%
}%


\ifwithsols
\title{حل امتحان مفمار في علم الحاسوب للصف الثامن - بايثون $2019$}
\else
\title{امتحان مفمار في علم الحاسوب للصف الثامن - بايثون $2019$}
\fi

\begin{document}

\maketitle
\thispagestyle{fancy}

\ifwithsols
\else
\begin{boxCode}
\begin{itemize}[nosep]
    \item \textbf{الفصل الأول (إجباري - $60$ علامة):} أجب عن \textbf{جميع} الأسئلة $1$-$5$.
    \item \textbf{الفصل الثاني (اختياري - $40$ علامة):} أجب عن \textbf{سؤالين} من بين الأسئلة $6$-$9$.
    \item \textbf{مدة الامتحان:} $120$ دقيقة.
    \item \textbf{مواد مساعدة مسموحة:} كل مادة مساعدة شخصية، باستثناء حاسبة أو حاسوب قابل للبرمجة.
    \item اقرأ أسئلة الامتحان بتمعّن وأجب عنها بانتباه في المكان المخصّص لذلك.
\end{itemize}
\end{boxCode}
\fi

\section*{الفصل الأول — إجباري ($60$ علامة)}
في هذا الفصل خمسة أسئلة. أجب عن \textbf{جميع} الأسئلة $1$-$5$.

\begin{enumerate}[label=\textbf{السؤال \arabic*:}, wide=0pt, leftmargin=*, itemsep=2.2em]

% ---------------------------------------------------------------
\item ($8$ علامات) أمامك جدول فيه ثلاثة مقاطع برنامج بلغة بايثون. لكل مقطع برنامج في الجدول، اكتب ما هو المُخرج الناتج في العمود الفارغ.

{\renewcommand{\arraystretch}{1.4}
\begin{longtable}{|C{1.5cm}|p{9cm}|C{4cm}|}
\hline
\textbf{البند} & \centering\textbf{كود بلغة بايثون} & \textbf{المُخرج الناتج} \\
\hline
$1$. \newline ($1$ علامة) &
\begin{pycell}
\begin{english}
\begin{minted}{python}
num1 = 2
num2 = num1 + 5
num1 = num1 + num2 + 3
print("num1 = ", num1, "num2 = ", num2)
\end{minted}
\end{english}
\end{pycell}
& \ifwithsols \begin{english}\small\texttt{num1 = 12 num2 = 7}\end{english} \fi \\
\hline
$2$. \newline ($1$ علامة) &
\begin{pycell}
\begin{english}
\begin{minted}{python}
num = 3
num = 10
num = num + 3
print("num = ", num)
\end{minted}
\end{english}
\end{pycell}
& \ifwithsols \begin{english}\small\texttt{num = 13}\end{english} \fi \\
\hline
$3$. \newline ($6$ علامات) &
\begin{pycell}
\begin{english}
\begin{minted}{python}
string = "Yes We Can!!!"
print(string[4])
print(len(string))
print(string[-3])
print(string[::2])
print(string[4:-3])
print(string[:])
\end{minted}
\end{english}
\end{pycell}
& \ifwithsols \begin{english}\small\texttt{W}\newline\texttt{13}\newline\texttt{!}\newline\texttt{YsW a!!}\newline\texttt{We Can}\newline\texttt{Yes We Can!!!}\end{english} \fi \\
\hline
\end{longtable}}
\filbreak

% ---------------------------------------------------------------
\item ($8$ علامات) أمامك جدول فيه أربعة مقاطع برنامج. لكل واحد من مقاطع البرنامج في الجدول، أعطِ مثالًا لمُدخل يعرض \textenglish{Yes} ومثالًا لمُدخل يعرض \textenglish{No}.

{\renewcommand{\arraystretch}{1.4}
\begin{longtable}{|C{1.4cm}|p{8.2cm}|C{2.4cm}|C{2.4cm}|}
\hline
\textbf{البند} & \centering\textbf{كود بلغة بايثون} & \textbf{مثال لمُدخل يعرض \textenglish{Yes}} & \textbf{مثال لمُدخل يعرض \textenglish{No}} \\
\hline
$1$. \newline ($2$ علامتان) &
\begin{pycell}
\begin{english}
\begin{minted}{python}
num1 = int(input("Enter first number"))
num2 = int(input("Enter second number"))
if num1 > 5 and num2 <= 4:
    print("Yes")
else:
    print("No")
\end{minted}
\end{english}
\end{pycell}
& \textenglish{num1:} \ifwithsols $\mathbf{6}$ \else \underline{\hspace{0.8cm}} \fi \newline \textenglish{num2:} \ifwithsols $\mathbf{4}$ \else \underline{\hspace{0.8cm}} \fi
& \textenglish{num1:} \ifwithsols $\mathbf{5}$ \else \underline{\hspace{0.8cm}} \fi \newline \textenglish{num2:} \ifwithsols $\mathbf{5}$ \else \underline{\hspace{0.8cm}} \fi \\
\hline
$2$. \newline ($2$ علامتان) &
\begin{pycell}
\begin{english}
\begin{minted}{python}
num1 = int(input("Enter first number"))
num2 = int(input("Enter second number"))
if num1 > 5 or num2 <= 4:
    print("Yes")
else:
    print("No")
\end{minted}
\end{english}
\end{pycell}
& \textenglish{num1:} \ifwithsols $\mathbf{6}$ \else \underline{\hspace{0.8cm}} \fi \newline \textenglish{num2:} \ifwithsols $\mathbf{10}$ \else \underline{\hspace{0.8cm}} \fi
& \textenglish{num1:} \ifwithsols $\mathbf{5}$ \else \underline{\hspace{0.8cm}} \fi \newline \textenglish{num2:} \ifwithsols $\mathbf{5}$ \else \underline{\hspace{0.8cm}} \fi \\
\hline
$3$. \newline ($2$ علامتان) &
\begin{pycell}
\begin{english}
\begin{minted}{python}
num1 = int(input("Enter first number"))
num2 = int(input("Enter second number"))
if (num1 > 5 or num2 <= 4) and num1 == num2:
    print("Yes")
else:
    print("No")
\end{minted}
\end{english}
\end{pycell}
& \textenglish{num1:} \ifwithsols $\mathbf{6}$ \else \underline{\hspace{0.8cm}} \fi \newline \textenglish{num2:} \ifwithsols $\mathbf{6}$ \else \underline{\hspace{0.8cm}} \fi
& \textenglish{num1:} \ifwithsols $\mathbf{6}$ \else \underline{\hspace{0.8cm}} \fi \newline \textenglish{num2:} \ifwithsols $\mathbf{4}$ \else \underline{\hspace{0.8cm}} \fi \\
\hline
$4$. \newline ($2$ علامتان) &
\begin{pycell}
\begin{english}
\begin{minted}{python}
num1 = int(input("Enter first number"))
num2 = int(input("Enter second number"))
if num1 > 5 or (num2 <= 4 and num1 == num2):
    print("Yes")
else:
    print("No")
\end{minted}
\end{english}
\end{pycell}
& \textenglish{num1:} \ifwithsols $\mathbf{7}$ \else \underline{\hspace{0.8cm}} \fi \newline \textenglish{num2:} \ifwithsols $\mathbf{10}$ \else \underline{\hspace{0.8cm}} \fi
& \textenglish{num1:} \ifwithsols $\mathbf{5}$ \else \underline{\hspace{0.8cm}} \fi \newline \textenglish{num2:} \ifwithsols $\mathbf{4}$ \else \underline{\hspace{0.8cm}} \fi \\
\hline
\end{longtable}}
\filbreak

% ---------------------------------------------------------------
\item ($15$ علامة) أمامك جدول فيه أربعة مقاطع برنامج. لكل مقطع برنامج في الجدول، اكتب ما هو المُخرج الناتج في العمود الفارغ.

{\renewcommand{\arraystretch}{1.4}
\begin{longtable}{|C{1.8cm}|p{8.5cm}|C{4cm}|}
\hline
\textbf{البند} & \centering\textbf{كود بلغة بايثون} & \textbf{المُخرج الناتج} \\
\hline
$1$. \newline ($3$ علامات) &
\begin{pycell}
\begin{english}
\begin{minted}{python}
for k in range(10):
    if k % 2 == 0:
        print(k + 1)
\end{minted}
\end{english}
\end{pycell}
& \ifwithsols \begin{english}\texttt{1}\newline\texttt{3}\newline\texttt{5}\newline\texttt{7}\newline\texttt{9}\end{english} \fi \\
\hline
$2$. \newline ($4$ علامات) &
\begin{pycell}
\begin{english}
\begin{minted}{python}
for k in range(20):
    if k % 3 == 0 and k <= 9:
        print(k)
\end{minted}
\end{english}
\end{pycell}
& \ifwithsols \begin{english}\texttt{0}\newline\texttt{3}\newline\texttt{6}\newline\texttt{9}\end{english} \fi \\
\hline
$3$. \newline ($4$ علامات) &
\begin{pycell}
\begin{english}
\begin{minted}{python}
for k in range(3, 10, 2):
    print(k)
\end{minted}
\end{english}
\end{pycell}
& \ifwithsols \begin{english}\texttt{3}\newline\texttt{5}\newline\texttt{7}\newline\texttt{9}\end{english} \fi \\
\hline
$4$. \newline ($4$ علامات) &
\begin{pycell}
\begin{english}
\begin{minted}{python}
for k in range(40, 0, -10):
    print(k)
\end{minted}
\end{english}
\end{pycell}
& \ifwithsols \begin{english}\texttt{40}\newline\texttt{30}\newline\texttt{20}\newline\texttt{10}\end{english} \fi \\
\hline
\end{longtable}}
\filbreak

% ---------------------------------------------------------------
\item ($14$ علامة) أمامك جدول فيه أربعة مقاطع برنامج. لكل مقطع برنامج في الجدول، اكتب ما هو المُخرج الناتج في العمود الفارغ.

{\renewcommand{\arraystretch}{1.4}
\begin{longtable}{|C{1.8cm}|p{8.5cm}|C{4cm}|}
\hline
\textbf{البند} & \centering\textbf{كود بلغة بايثون} & \textbf{المُخرج الناتج} \\
\hline
$1$. \newline ($3$ علامات) &
\begin{pycell}
\begin{english}
\begin{minted}{python}
number = 123
print(number // 100)
print(number // 10 % 10)
print(number % 10)
\end{minted}
\end{english}
\end{pycell}
& \ifwithsols \begin{english}\texttt{1}\newline\texttt{2}\newline\texttt{3}\end{english} \fi \\
\hline
$2$. \newline ($2$ علامتان) &
\begin{pycell}
\begin{english}
\begin{minted}{python}
number = 12345
print(number // 100)
print(number % 100)
\end{minted}
\end{english}
\end{pycell}
& \ifwithsols \begin{english}\texttt{123}\newline\texttt{45}\end{english} \fi \\
\hline
$3$. \newline ($2$ علامتان) &
\begin{pycell}
\begin{english}
\begin{minted}{python}
number = 12345
print(number // 100 // 10)
print(number // 10 % 10)
\end{minted}
\end{english}
\end{pycell}
& \ifwithsols \begin{english}\texttt{12}\newline\texttt{4}\end{english} \fi \\
\hline
$4$. \newline ($7$ علامات) &
\begin{pycell}
\begin{english}
\begin{minted}{python}
number = 123
while number > 0:
    print(number % 10)
    number = int(number // 10)
\end{minted}
\end{english}
\end{pycell}
& \ifwithsols \begin{english}\texttt{3}\newline\texttt{2}\newline\texttt{1}\end{english} \fi \\
\hline
\end{longtable}}
\filbreak

% ---------------------------------------------------------------
\item ($15$ علامة) أمامك ثلاثة مقاطع برنامج تتضمّن عملية باسم \textenglish{num\_even}.
الهدف من كل واحد من مقاطع البرنامج هو استقبال $10$ أعداد صحيحة، وطباعة كمية الأعداد \textbf{ثنائية المنازل الموجبة} (من $10$ إلى $99$) التي ظهرت في المُدخل.
\textbf{مثال:} لسلسلة الأعداد العشرة التالية: $230$، $-555$، $76$، $5$، $80$، $-100$، $77$، $8$، $74$، $90$، يُطبع العدد $5$.
في كل مقطع برنامج وقع خطأ منطقي واحد، ولذلك لا يحقّق المقطع هدفه. اشرح الخطأ وصحّحه.
(ملاحظة: الخطأ ليس خطأ وقت تشغيل. افترض أنّ المُدخل سليم ويشمل أعدادًا صحيحة فقط.)

\begin{enumerate}[label=\textbf{مقطع \alph*:}, wide=0pt, leftmargin=*, itemsep=1.5em]

\item .
\begin{boxCode}
\begin{english}
\begin{minted}{python}
def num_even():
    count = 0
    index = 0
    num = int(input("Enter an integer number"))
    while index < 10:
        if num <= 99 and num >= 10:
            count = count + 1
        index = index + 1
    print("count", count)
\end{minted}
\end{english}
\end{boxCode}
\ifwithsols
\begin{boxSolution}
\textbf{الخطأ:} أمر الاستقبال \textenglish{\texttt{num = int(input(...))}} مكتوب خارج الحلقة فقط، ولذلك يستقبل البرنامج عددًا واحدًا فقط ويفحصه $10$ مرات بدل استقبال $10$ أعداد مختلفة. \\
\textbf{التصحيح:} نقل (أو إضافة) السطر \textenglish{\texttt{num = int(input("Enter an integer number"))}} إلى داخل حلقة \textenglish{while}.
\end{boxSolution}
\else
\textbf{ما هو الخطأ؟} \underline{\hspace{11cm}} \\
\textbf{كيف يجب تصحيحه؟} \underline{\hspace{10cm}}
\fi
\filbreak

\item .
\begin{boxCode}
\begin{english}
\begin{minted}{python}
def num_even():
    count = 0
    index = 0
    while index < 10:
        num = int(input("Enter an integer number"))
        if num <= 99 and num >= 10:
            print("count", count)
        index = index + 1
\end{minted}
\end{english}
\end{boxCode}
\ifwithsols
\begin{boxSolution}
\textbf{الخطأ:} لا يتمّ زيادة العداد \textenglish{count} بـ $1$ عند تحقق الشرط، وبدلًا من ذلك تتمّ طباعة \textenglish{count} (الذي يبقى $0$) داخل الشرط بدل طباعته في نهاية العملية بعد انتهاء الحلقة. \\
\textbf{التصحيح:} استبدال أمر الطباعة داخل الشرط بـ \textenglish{\texttt{count = count + 1}}، ووضع \textenglish{\texttt{print("count", count)}} خارج حلقة \textenglish{while} في نهاية العملية.
\end{boxSolution}
\else
\textbf{ما هو الخطأ؟} \underline{\hspace{11cm}} \\
\textbf{كيف يجب تصحيحه؟} \underline{\hspace{10cm}}
\fi
\filbreak

\item .
\begin{boxCode}
\begin{english}
\begin{minted}{python}
def num_even():
    count = 0
    index = 0
    while index < 10:
        if num <= 99 and num >= 10:
            count = count + 1
            num = int(input("Enter an integer number"))
        index = index + 1
    print("count", count)
\end{minted}
\end{english}
\end{boxCode}
\ifwithsols
\begin{boxSolution}
\textbf{الخطأ:} أمر الاستقبال \textenglish{\texttt{num = int(input(...))}} موضوع داخل جملة الشرط \textenglish{if} (وبعد الفحص)، بدل أن يكون في بداية كل تكرار من الحلقة قبل جملة \textenglish{if}. \\
\textbf{التصحيح:} نقل السطر \textenglish{\texttt{num = int(input("Enter an integer number"))}} ليكون السطر الأول داخل حلقة \textenglish{while} قبل جملة الشرط \textenglish{if}.
\end{boxSolution}
\else
\textbf{ما هو الخطأ؟} \underline{\hspace{11cm}} \\
\textbf{كيف يجب تصحيحه؟} \underline{\hspace{10cm}}
\fi
\filbreak

\end{enumerate}
\filbreak

\end{enumerate}

\clearpage
\section*{الفصل الثاني — اختياري ($40$ علامة)}
أجب عن \textbf{سؤالين} من بين الأسئلة $6$-$9$ (لكل سؤال — $20$ علامة).

\begin{enumerate}[resume, label=\textbf{السؤال \arabic*:}, wide=0pt, leftmargin=*, itemsep=2.2em]

% ---------------------------------------------------------------
\item ($20$ علامة) أمامك مقطع برنامج يستعمل مكتبة \textenglish{turtle}:

\begin{boxCode}
\begin{english}
\begin{minted}{python}
import turtle

tw = turtle.Screen()
player = turtle.Turtle()
player.shape("square")
size = 100
player.pendown()

for index in range(4):
    for item in range(4):
        player.forward(size)
        player.left(90)
    size = size - 20

turtle.mainloop()
\end{minted}
\end{english}
\end{boxCode}

تمّ تشغيل مقطع البرنامج. أيّ من المخرجات التالية سيُعرض على شاشة الحاسوب؟ أحط بدائرة الإجابة الصحيحة:

{\centering
\begin{english}
\renewcommand{\arraystretch}{1.2}
\begin{tabular}{|C{6.5cm}|C{6.5cm}|}
\hline
\vspace{0.2cm}
\begin{tikzpicture}[scale=0.6]
\draw[thick] (0,0) rectangle (4,1);
\draw[thick] (1,0) -- (1,1);
\draw[thick] (2,0) -- (2,1);
\draw[thick] (3,0) -- (3,1);
\draw[-{Latex[length=2mm]}, thick] (3.5,0) -- (4,0);
\end{tikzpicture}
\newline \textarabic{\textbf{أ}}
&
\vspace{0.2cm}
\begin{tikzpicture}[scale=0.6]
\draw[thick] (0,0) rectangle (3.2,3.2);
\draw[thick] (0,0) rectangle (2.5,2.5);
\draw[thick] (0,0) rectangle (1.8,1.8);
\draw[thick] (0,0) rectangle (1.1,1.1);
\draw[-{Latex[length=2mm]}, thick] (-0.8,0) -- (0,0);
\end{tikzpicture}
\newline \textarabic{\textbf{ب}}
\\ \hline
\vspace{0.2cm}
\begin{tikzpicture}[scale=0.6]
\draw[thick] (0,0) rectangle (3.2,3.2);
\draw[thick] (0.4,0.4) rectangle (2.8,2.8);
\draw[thick] (0.8,0.8) rectangle (2.4,2.4);
\draw[thick] (1.2,1.2) rectangle (2.0,2.0);
\draw[-{Latex[length=2mm]}, thick] (0,2.5) -- (0,3.2);
\end{tikzpicture}
\newline \textarabic{\textbf{ج}}
&
\vspace{0.2cm}
\begin{tikzpicture}[scale=0.6]
\draw[thick] (0,0) rectangle (1,4);
\draw[thick] (0,1) -- (1,1);
\draw[thick] (0,2) -- (1,2);
\draw[thick] (0,3) -- (1,3);
\draw[-{Latex[length=2mm]}, thick] (0,3.5) -- (0,4);
\end{tikzpicture}
\newline \textarabic{\textbf{د}}
\\ \hline
\end{tabular}
\end{english}
\par} %

\ifwithsols
\begin{boxSolution}
الإجابة الصحيحة هي \textbf{ب}. \\
\textbf{التعليل:} الحلقة الداخلية ترسم مربّعًا طول ضلعه \textenglish{size} وتعود لنفس نقطة البداية وبنفس الاتجاه (لأنّها تدور $90^\circ$ يسارًا $4$ مرات، أي $360^\circ$). في كل تكرار للحلقة الخارجية يقلّ طول الضلع \textenglish{size} بمقدار $20$، فترسم $4$ مربعات متداخلة تشترك جميعها في الزاوية اليسرى السفلية $(0,0)$.
\end{boxSolution}
\fi
\filbreak

% ---------------------------------------------------------------
\item ($20$ علامة) متوالية فيبوناتشي (\textenglish{Fibonacci}) هي متوالية يكون حدّاها الأولان $1$، وكل حدّ يليهما يساوي مجموع الحدّين اللذين سبقاه.
أمامك متوالية فيبوناتشي المكوّنة من $8$ حدود (من اليسار إلى اليمين):
$$1, 1, 2, 3, 5, 8, 13, 21$$
معروف أنّ الحدّ الأول في المتوالية هو $\text{first} = 1$، والحدّ الثاني هو $\text{second} = 1$.

\begin{enumerate}[label=\alph*.]
\item ($2$ علامتان) ما هو الحدّ التالي في المتوالية (الحدّ التاسع) في المثال؟ \ifwithsols $\mathbf{34}$ (لأنّ $13 + 21 = 34$) \else \underline{\hspace{3cm}} \fi
\filbreak

\item ($18$ علامة) أمامك مقطع برنامج يطبع الحدود الـ $8$ الأولى في المتوالية. أكمل المقطع:

\begin{boxCode}
\begin{english}
\begin{minted}{python}
first = ____________
second = ____________
counter = 3
print("Fibonacci item 1 is ", first)
print("Fibonacci item 2 is ", second)
while counter <= 8:
    third = ______ + _______
    first = ____________
    second = ____________
    print("Fibonacci item ", counter, "is", ____________)
    counter = counter + ____________
\end{minted}
\end{english}
\end{boxCode}

\ifwithsols
\begin{boxSolution}
\begin{english}
\begin{minted}{python}
first = 1
second = 1
counter = 3
print("Fibonacci item 1 is ", first)
print("Fibonacci item 2 is ", second)
while counter <= 8:
    third = first + second
    first = second
    second = third
    print("Fibonacci item ", counter, "is", third)
    counter = counter + 1
\end{minted}
\end{english}
\end{boxSolution}
\fi
\filbreak
\end{enumerate}
\filbreak

% ---------------------------------------------------------------
\item ($20$ علامة) طبّقت إدارة مدرسة "علم وتكنولوجيا للجميع" برنامجًا تعليميًا جديدًا لطلاب طبقة الصف الثامن، والتي تضمّ $300$ طالب. لفحص ما إذا كان البرنامج قد ساهم في تحسين تحصيلات الطلاب، تقرّر المقارنة بين علامات الطلاب في نهاية النصف الأول من السنة (\textenglish{gradeA}) وبين علاماتهم في نهاية النصف الثاني (\textenglish{gradeB}).
من أجل الفحص، كُتب برنامج يستقبل لكل طالب العلامة التي حصل عليها في النصف الأول والعلامة التي حصل عليها في النصف الثاني. يعدّ البرنامج ويطبع كمية الطلاب الذين حسّنوا علاماتهم:
\begin{itemize}
\item يُعرَّف التحسّن فقط عندما تكون العلامة في النصف الثاني \textbf{أعلى من أو تساوي} العلامة في النصف الأول.
\item إذا حسّن أكثر من $200$ طالب علاماتهم في هذا الفحص، يُطبع \textenglish{"The program is successful"}؛ وإلّا يُطبع \textenglish{"The program needs improvement"}.
\end{itemize}
أمامك كود البرنامج. أكمل التعليمات الناقصة فيه (انتبه: للمتغيرات ووظيفتها في البرنامج أهمية):

\begin{boxCode}
\begin{english}
\begin{minted}{python}
def Exs8():
    counter = _______
    for k in range(_______):
        gradeA = int(input("Enter gradeA"))
        gradeB = int(input("Enter gradeB"))
        if gradeB >= gradeA:
            counter = _______ + _______
    if counter > _______:
        print("___________________________")
    else:
        print("___________________________")

Exs8()
\end{minted}
\end{english}
\end{boxCode}

\ifwithsols
\begin{boxSolution}
\begin{english}
\begin{minted}{python}
def Exs8():
    counter = 0
    for k in range(300):
        gradeA = int(input("Enter gradeA"))
        gradeB = int(input("Enter gradeB"))
        if gradeB >= gradeA:
            counter = counter + 1
    print(counter)
    if counter > 200:
        print("The program is successful")
    else:
        print("The program needs improvement")

Exs8()
\end{minted}
\end{english}
\end{boxSolution}
\fi
\filbreak

% ---------------------------------------------------------------
\item ($20$ علامة) كتبت عدي سيناريو بحيث تؤدّي كل نقرة (\textenglish{click}) على الفأرة إلى رسم مربّع على الشاشة في مكان النقرة:
\begin{itemize}
\item يُرسم المربّع على الشاشة باللونين الأحمر والأزرق بالتناوب (النقرة الأولى — أحمر \textenglish{"red"}، النقرة الثانية — أزرق \textenglish{"blue"}، وهكذا).
\item في كل نقرة يُرسم مربّع يزداد طول ضلعه بـ $2$.
\end{itemize}
أمامك مقطع برنامج تنقصه تعليمات. أكمل التعليمات الناقصة:
(ملاحظة: المتغير \textenglish{size} يمثّل طول ضلع المربّع، والمتغير \textenglish{click\_num} يمثّل رقم النقرة.)

\begin{boxCode}
\begin{english}
\begin{minted}{python}
import turtle

def do_it(x, y):
    global click_num
    global size
    player.penup()
    player.goto(x, y)
    player.shapesize(size)
    if ____________ % 2 == 0:
        player.color("_____")
    else:
        player.color("_____")
    player.stamp()
    click_num = ____________ + ____________
    size = ____________ + ____________

tw = turtle.Screen()
click_num = 0
size = 2
tw.onclick(do_it)
player = turtle.Turtle()
player.shape("square")
turtle.mainloop()
\end{minted}
\end{english}
\end{boxCode}

\ifwithsols
\begin{boxSolution}
\begin{english}
\begin{minted}{python}
import turtle

def do_it(x, y):
    global click_num
    global size
    player.penup()
    player.goto(x, y)
    player.shapesize(size)
    if click_num % 2 == 0:
        player.color("red")
    else:
        player.color("blue")
    player.stamp()
    click_num = click_num + 1
    size = size + 2

tw = turtle.Screen()
click_num = 0
size = 2
tw.onclick(do_it)
player = turtle.Turtle()
player.shape("square")
turtle.mainloop()
\end{minted}
\end{english}
\end{boxSolution}
\fi
\filbreak

\end{enumerate}

\vspace{1cm}
\begin{flushleft}
أرجو لكم وقتًا ممتعًا.

الأستاذ محمود اغبارية.
\end{flushleft}

\end{document}
